\documentclass[uplatex,a4paper,12pt]{jsarticle}

% レイアウト
\usepackage[margin=25mm]{geometry}
\usepackage{setspace}
\onehalfspacing

% 図表
\usepackage[dvipdfmx]{graphicx}
\usepackage{booktabs}

% リンク
\usepackage[dvipdfmx]{hyperref}
\usepackage{pxjahyper}
\hypersetup{
  colorlinks=true,
  linkcolor=blue,
  urlcolor=blue,
  citecolor=blue
}

% レポート情報
\title{テーマ}
%\author{学籍番号: 12345678 \\
%氏名: 山田 太郎}
\author{学籍番号: \texttt{6930\mbox{-}38\mbox{-}5047}\\
名前：程 ヨウ}

\date{\today}

\begin{document}

\maketitle
%\tableofcontents
%\newpage

\section{はじめに}
本レポートでは、研究の背景と目的を述べる。

\section{方法}
\subsection{実験条件}
ここに実験条件を記述する。

\subsection{データ収集}
ここにデータ収集方法を記述する。

\section{結果}
\subsection{図の例}
図\ref{fig:example}に結果の例を示す。

\begin{figure}[htbp]
  \centering
  % \includegraphics[width=0.7\textwidth]{example.png}
  \fbox{\rule{0pt}{40mm}\rule{0.7\textwidth}{0pt}}
  \caption{結果図の例}
  \label{fig:example}
\end{figure}

\subsection{表の例}
表\ref{tab:example}に結果の要約を示す。

\begin{table}[htbp]
  \centering
  \caption{結果の要約}
  \label{tab:example}
  \begin{tabular}{lcc}
    \toprule
    項目 & 条件A & 条件B \\
    \midrule
    精度 & 0.91 & 0.88 \\
    再現率 & 0.89 & 0.85 \\
    \bottomrule
  \end{tabular}
\end{table}

\section{考察}
結果から分かる点と課題を述べる。

\section{結論}
本レポートの結論を簡潔に述べる。

\begin{thebibliography}{9}
\bibitem{ref1}
著者名, 『文献タイトル』, 出版社, 年.
\end{thebibliography}

\end{document}
